\documentclass[9pt]{extarticle}
\usepackage[T1]{fontenc}
\usepackage[utf8]{inputenc}
\usepackage[french]{babel}
\usepackage{newpxtext,newpxmath}
\usepackage[a4paper,top=0.7cm,bottom=0.6cm,left=0.85cm,right=0.85cm]{geometry}
\usepackage{microtype}
\usepackage{xcolor}
\usepackage{titlesec}
\usepackage[most]{tcolorbox}
\usepackage{fancyvrb}
\usepackage{multicol}
\usepackage{enumitem}

\definecolor{navy}{HTML}{1F3A5F}
\definecolor{grey}{HTML}{54606E}
\definecolor{green}{HTML}{2C6E49}
\definecolor{red}{HTML}{A6322E}
\definecolor{amber}{HTML}{B25A1E}
\definecolor{boxgray}{HTML}{F3F5F8}
\definecolor{boxred}{HTML}{FCF2F1}
\definecolor{boxamber}{HTML}{FBF3EA}
\definecolor{boxedge}{HTML}{D9E0EA}

\setlength{\parindent}{0pt}
\setlength{\parskip}{1.2pt}
\setlength{\columnsep}{11pt}
\setlength{\emergencystretch}{3em}
\linespread{0.94}
\pagestyle{empty}

\titleformat{\section}{\color{navy}\bfseries}{}{0pt}{}
\titlespacing{\section}{0pt}{4.5pt}{1pt}

\newtcolorbox{these}{colback=boxgray,colframe=boxedge,boxrule=0.5pt,arc=2pt,
  left=6pt,right=6pt,top=3pt,bottom=3pt}
\newtcolorbox{alerte}{colback=boxred,colframe=red!55,boxrule=0.7pt,arc=2pt,
  left=6pt,right=6pt,top=3pt,bottom=3pt}
\newtcolorbox{cadre}{colback=boxamber,colframe=amber!55,boxrule=0.6pt,arc=2pt,
  left=6pt,right=6pt,top=3pt,bottom=3pt}
\newcommand{\code}[1]{\texttt{\footnotesize #1}}
\newcommand{\mm}[1]{\textbf{\textcolor{navy}{#1}}}
\newcommand{\ok}[1]{\textcolor{green}{\textbf{#1}}}
\newcommand{\ko}[1]{\textcolor{red}{\textbf{#1}}}
\DefineVerbatimEnvironment{cmd}{Verbatim}{fontsize=\footnotesize,xleftmargin=6pt,
  commandchars=\\\{\}}

\begin{document}

\begin{center}
{\LARGE\bfseries\color{navy} Le Seuil --- une commande, et ce qu'elle prouve}\\[2pt]
{\small\color{grey}\itshape La page à lire \emph{avant} de la lancer : ce qui va défiler, ce que vous
pourrez recompter vous-même, et ce que cela ne prouvera pas}\\[2pt]
{\footnotesize\color{grey} 4 septembre 2026 \; · \; chaque chiffre de cette page a été mesuré ce
jour-là, sur cette démonstration, et non repris d'un document antérieur}
\end{center}
\vspace{1pt}

\begin{these}
Une passerelle qui tourne \textbf{sur votre machine} sert un reçu, l'exporte, en publie
l'empreinte \emph{hors} du reçu, puis en coupe une copie --- et la coupe se voit \textbf{par un
comptage}, sans clé, sans cryptographie, sans avoir à croire qui que ce soit. C'est tout ce que
la démonstration prétend. Ce qu'elle ne prétend pas est écrit au bas de cette page, avec la même
précision.
\end{these}

\section{La commande}
\begin{cmd}
docker compose up
\end{cmd}
Rien d'autre. Aucun compte, aucune clé d'API, aucune inscription. Deux services démarrent : la
\mm{passerelle}, qui reste vivante pour que vous puissiez l'interroger, et \mm{preuve}, un pilote
qui conduit la démonstration puis s'arrête. \code{Ctrl-C} pour sortir.

\begin{cadre}
\textbf{Si rien ne démarre, c'est presque toujours le port.} \code{8080} est occupé sur beaucoup de
machines --- il l'était sur celle-ci le jour de la mesure. Le port hôte se déplace sans éditer
un fichier :\; \code{SCELLA\_PORT=18080 docker compose up}\;. La passerelle n'est publiée que
sur la boucle locale \code{127.0.0.1}, jamais sur toutes les interfaces.
\end{cadre}

\begin{multicols}{2}

\section{Ce qui défile}
Le pilote conduit \textbf{cinq étapes} et les numérote. Elles disent, dans l'ordre :

\begin{enumerate}[leftmargin=1.15em,itemsep=1pt,topsep=1.5pt]
\item la passerelle répond --- et le rapport de démarrage annonce
      «~provenance \textbf{DÉGRADÉE}~», parce qu'aucune clé n'a été fournie. Elle le
      \emph{dit} ; elle ne le tait pas ;
\item un run réel est servi, avec son identifiant et sa racine ;
\item le reçu est exporté : \code{bundle.sealed.jsonl}, \textbf{3 nœuds} --- et son \textbf{ancre}
      est publiée \emph{hors} du fichier ;
\item une copie est tronquée. L'ancre annonce \textbf{3}, l'intact en contient \textbf{3}, le
      tronqué \textbf{2} : la contradiction se lit sans rien déchiffrer ;
\item ce qui a été prouvé, et ce qui ne l'a pas été --- la démonstration se clôt en
      énonçant ses propres limites.
\end{enumerate}

\section{Ce que vous pouvez recompter}
Le volume \code{scella\_preuve} porte \textbf{trois fichiers}, et rien n'oblige à croire le pilote
sur parole. Sortez-les :
\begin{cmd}
docker compose cp preuve:/preuve ./p
\end{cmd}
\vspace{-2pt}
\code{anchors.log} (157~o) · \code{bundle.sealed.jsonl} (1\,340~o) ·
\code{bundle.TRONQUE.jsonl} (893~o).

\smallskip
L'ancre est une ligne à \textbf{six champs tabulés} : \code{\$2} le run, \code{\$4} la racine,
\code{\$5} le nombre de n\oe uds. Le comptage, refait à la main :

\begin{cmd}
awk -F'\\t' -v r=<run> '$2==r \{print $5\}' \\
    p/anchors.log            \textcolor{green}{→ 3}
grep -c '[^[:space:]]' \\
    p/bundle.sealed.jsonl    \textcolor{green}{→ 3}
grep -c '[^[:space:]]' \\
    p/bundle.TRONQUE.jsonl   \textcolor{red}{→ 2}
\end{cmd}

\textbf{Par le run, jamais par la position.} Sous le profil \code{jonction}, deux passerelles
publient dans ce journal : «~la dernière ligne~» n'est pas forcément la vôtre. Le run se lit dans
le bundle lui-même --- c'est le préfixe de \code{node\_id}, avant le \code{::}.

\section{Le détail qui explique tout}
La racine annoncée par l'ancre \textbf{est le sceau du dernier maillon}. Elle est donc dans le
bundle intact, et elle \emph{disparaît} avec la coupe :

\begin{cmd}
grep -c <racine> p/bundle.sealed.jsonl  \textcolor{green}{→ 1}
grep -c <racine> p/bundle.TRONQUE.jsonl \textcolor{red}{→ 0}
\end{cmd}

Le témoin part avec le crime. C'est pourquoi un bundle tronqué \textbf{se re-prouve
parfaitement} --- un préfixe de chaîne valide reste valide --- et pourquoi une empreinte reçue
par un canal \emph{distinct} est la seule chose qui puisse le contredire. L'image porte deux
vérifieurs : \code{verify-mini} (1\,159 lignes, \textbf{arbre de dépendances vide} : un auditeur
peut le lire en entier) et \code{provenance-verify}, le seul des deux qui sache lire un journal
d'ancres (\code{-{}-anchor}).

\section{L'image}
\code{scella-coin:local}, \textbf{54,2~Mo}, \textbf{distroless} : ni shell, ni gestionnaire de
paquets, ni utilitaire. Quatre binaires, rien d'autre. C'est ce qu'un inconnu exécute en premier
--- le pire moment pour lui offrir une surface d'attaque. C'est aussi pourquoi le pilote est un
binaire et non un script, et pourquoi on lit le volume avec \code{docker compose cp} (exécuté par
le démon) plutôt qu'avec un \code{ls} qui n'existe pas dans l'image.

\smallskip
\textbf{Aucun octet ne sort à l'exécution}, et ce n'est pas une promesse : le chemin d'appel vers
un fournisseur tiers vit sous \code{\#[cfg(feature = "frontier-llm")]}, et l'image du Seuil est
construite \emph{sans aucune feature}. Le code qui appellerait dehors n'est pas compilé.

\end{multicols}

\begin{alerte}
\textbf{CE QUE CETTE DÉMONSTRATION NE PROUVE PAS.} Quatre choses, et il vaut mieux les lire ici
que les découvrir plus tard.

\smallskip
\textbf{1. Elle ne prouve rien \emph{contre} quelqu'un.} Aucune clé n'est injectée : la passerelle
tourne sur la clé de \textbf{test publique}, connue de quiconque a lu le dépôt. Le reçu démontre
le \textbf{mécanisme}, jamais la \textbf{substance}. Pour qu'il engage, il faut une clé que vous
seul détenez (\code{EFE\_NANO\_AUDIT\_HMAC\_MASTER}, injectée au déploiement) --- et alors la même
commande produit un reçu opposable.

\textbf{2. Le \emph{build}, lui, télécharge.} La promesse «~aucun octet qui sort~» porte sur
l'\textbf{exécution}. La construction de l'image tire les images de base et les caisses depuis le
réseau. Sur une machine cloisonnée, c'est là que l'on bute --- et pas sur le produit.

\textbf{3. Un reçu atteste l'\emph{intégrité} de ce qui a été enregistré, jamais la
\emph{correction} de la décision.} Il dit «~voici exactement ce qui s'est passé, et personne ne
l'a réécrit depuis~». Il ne dit pas que ce qui s'est passé était juste. Aucune quantité de
cryptographie ne comble cet écart.

\textbf{4. «~Machine vierge~» n'a pas été mesuré.} Ce qui a été mesuré, c'est le comportement sur
une machine qui portait déjà un état Docker considérable. C'est une information utile ; ce n'est
pas celle que la formule annonce.
\end{alerte}

\vspace{1.5pt}
{\footnotesize\color{grey}
\textbf{Et ensuite ?} \code{docker compose -{}-profile jonction up} ajoute \mm{Le Registre} : le
\emph{même} run produit en plus un dossier réglementaire, nommé par l'identifiant du run, avec sa
vue de cadre, son mapping daté et son mode d'emploi de re-dérivation. Ce dossier ne prononce pas
la conformité --- outillé n'est pas conforme.
\hfill\textbf{Contrôle :} \code{bash scripts/e1-check.sh} rejoue tout ceci et recompte lui-même,
sans faire confiance au pilote.}

\end{document}
